\section{Experimental Evaluation}\label{s:experiments}

We developed a Layer-1 implementation with CATX integrated as described above. 
It combines \texttt{malachite} BFT consensus engine~\cite{malachite_repo} with Ethereum's \texttt{reth} execution client~\cite{reth_repo}.
It provides EVM compatibility with CATX supporting ECDSA, Falcon-512, and ML-DSA-44.
All experiments and benchmarks were run in a local 3 node network on an Apple M3 Max with 36 GiB of memory. 
We also measured SLH-DSA~\cite{NIST:24SLH} variants (SHAKE-128f and SHAKE-128s), but report them only in microbenchmarks. In end-to-end experiments, SHAKE-128f began to degrade around 200~TPS and failed at 300~TPS, while SHAKE-128s sustained only about 16 TPS and proved impractical to benchmark locally due to its 8{,}576$\times$ higher signing cost compared to ECDSA.

In this section, we evaluate the impact of post-quantum signatures on the blockchain performance. 
Our evaluation spans over 30{,}000 blocks and 11 million transactions under sustained load, enabling analysis of steady-state behavior. 
Under these conditions, throughput is limited by three resources: (1)~\emph{CPU}, when verification time approaches the 12\,s slot limit~\cite{ETH:W14}; (2)~\emph{gas}, when the block gas limit is reached; and (3)~\emph{block size}, when the block size cap is hit. 
Identifying the active bottleneck is essential to explain performance differences across signature schemes. 
We begin with microbenchmarks of the cryptographic primitives to isolate the effects of computation, transaction size, and network propagation on overall performance.


\subsection{Signatures Microbenchmarks}\label{subsec:crypto-primitives}

We start by establishing a computational baseline through isolated benchmarks.
Table~\ref{tab:crypto-ops} reports key generation, signing, and verification times for ECDSA and post-quantum schemes Falcon-512, ML-DSA-44 and SLH-DSA for SHAKE-128f and SHAKE-128s variants.
% \vspace{-20pt}
% \begin{table}[htbp]
% \centering
% \caption{Median of cryptographic operation performance ($\mu$s) using Rust crate \texttt{criterion}~\cite{criterion_rs} with 100 samples and 5-second measurement periods.}
% \label{tab:crypto-ops}
% \setlength{\tabcolsep}{8pt}
% \begin{tabular}{@{}lrrr@{}}
% \toprule
% \textbf{Algorithm} & \textbf{Key gen.} & \textbf{Signing} & \textbf{Verification}  \\
% \midrule
% ECDSA & 34.3 & 49.1 & 54.7  \\
% Falcon-512 & 4,386 & 127.7 & \textbf{19.9} \\
% ML-DSA-44 & 40.0 & 179.3 & \textbf{43.1} \\
% \bottomrule
% \end{tabular}
% \end{table}



\begin{table}[htbp]
    \centering
    \caption{Public key and signature sizes (bytes) and median cryptographic operation performance ($\mu$s) using Rust crate \texttt{criterion}~\cite{criterion_rs} with 100 samples and 5-second measurement periods. Values in parentheses are relative to ECDSA.}
    \label{tab:crypto-ops}
    \setlength{\tabcolsep}{8pt}
    \resizebox{\linewidth}{!}{%
    \begin{tabular}{@{}lrrr@{\hspace{10pt}}rrr@{}}
    \toprule
    & \multicolumn{3}{c}{\textbf{Sizes (B)}} & \multicolumn{3}{c}{\textbf{Latency $\mu$s}} \\
    \cmidrule(lr){2-4}\cmidrule(lr){5-7}
    \textbf{Scheme} & \textbf{Public Key} & \textbf{Signature} & \textbf{Total} & \textbf{Key Gen.} & \textbf{Sign} & \textbf{Verify} \\
    \midrule
    ECDSA & 33 & 65 & 98 & 31 & 45 & 51 \\
    Falcon-512 & 897 & 666 & 1{,}563 (16.0$\times$) & 4{,}035 & 119\,(2.6$\times$) & \textbf{18\,(0.4$\times$)} \\
    ML-DSA-44 & 1{,}312 & 2{,}420 & 3{,}732 (38.1$\times$) & 37 & 165\,(3.7$\times$) & \textbf{40\,(0.8$\times$)} \\
    SLH-DSA-128f & 32 & 17{,}088 & 17{,}120 (175$\times$) & 808 & 19{,}185\,(426$\times$) & 1{,}154\,(23$\times$) \\
    SLH-DSA-128s & 32 & 7{,}856 & 7{,}888 (80.5$\times$) & 49{,}985 & 385{,}903\,(8{,}576$\times$) & 371\,(7.3$\times$) \\
    \bottomrule
    \end{tabular}
    }
\end{table}%


% \vspace{-10pt}
We observe that for both Falcon-512 and ML-DSA-44 schemes, signature verification -- the dominant cryptographic operation executed by the validating nodes -- is faster than ECDSA~\cite{ESORICS:KARR25}.
% Falcon-512 achieves a 2.75$\times$ reduction relative to ECDSA, while ML-DSA-44 achieves a 1.27$\times$ reduction.
These measurements indicate that post-quantum verification does not introduce a computational bottleneck for these two schemes.
On the other hand, both SLH-DSA variants exceed the verification latency between 7$\times$  and 23$\times$.

Signing costs are higher for all post-quantum schemes, where the SLH-DSA variant optimized for size being the most affected (8{,}576$\times$).
However, signing occurs off-chain and, therefore, does not affect consensus throughput.
In addition, key generation is a one-time account initialization cost, which is amortized over its lifetime without affecting performance.

In contrast, signature and public key sizes differ by more than two orders of magnitude across schemes and have a direct effect on transaction size and block utilization.
Table~\ref{tab:crypto-ops} summarizes these size differences as well.
% % \vspace{-2pt}
% \begin{table}[htbp]
% \centering
% \caption{Signature and public key sizes in bytes B.}
% \label{tab:sig-sizes}
% \setlength{\tabcolsep}{8pt}
% \begin{tabular}{@{}lrrrl@{}}
% \toprule
% \textbf{Algorithm} & \textbf{Public key} & \textbf{Signature} & \textbf{Total} & \textbf{Overhead} \\
% \midrule
% ECDSA & 65~B & 65~B & 130~B & 1.0$\times$ \\
% Falcon-512 & 897~B & 666~B & 1,563~B & \textbf{12.0$\times$} \\
% ML-DSA-44 & 1,312~B & 2,420~B & 3,732~B & \textbf{28.7$\times$} \\
% \bottomrule
% \end{tabular}
% \end{table}
% \vspace{-10pt}

\subsection{Gas Computation}

As mentioned before, gas costs have an impact on blockchain throughput. 
These are computed from three components: (1)~intrinsic gas, a base cost of 21,000~gas per transaction; (2)~data gas, charged at 4~gas per zero byte~\cite{ETH:W14} and 16~gas per non-zero byte~\cite{EIP2028} in transaction data; and (3)~execution gas, for operations like signature verification via precompiles.

For signature verification, ECDSA uses a fixed 3,000~gas precompile cost~\cite{ETH:W14}, while post-quantum schemes use base costs plus per-message length charges: Falcon-512 requires 1,300~gas base plus 12~gas per 32-byte word of message data, and ML-DSA-44 requires 2,800~gas base plus 12~gas per word.
These execution gas parameters are derived by scaling proportionally to measured signature verification times, using ECDSA as the reference.%, ensuring gas reflects actual computational work.

\subsection{CATX Impact}\label{subsec:tx-impact}

We next evaluate CATX transactions (cf. Section~\ref{s:catx}) under sustained load to quantify end-to-end throughput. 

We first investigate the transaction sizes and the gas threshold at which blocks become constrained by the 8~MB block size rather than by the gas or CPU limit, reported in Table~\ref{tab:tx-sizes}.
\begin{table}[htbp]
\centering
\caption{Transaction sizes and gas at 8 MB block size limit. B: billion; M: million.}
\label{tab:tx-sizes}
\setlength{\tabcolsep}{8pt}
\begin{tabular}{@{}lrrl@{}}
\toprule
\textbf{Type} & \textbf{TX size} & \textbf{Data-limited at} & \textbf{Practical impact} \\
\midrule
Legacy\footnotemark & 121~B & 1.456~B gas & Never data-limited \\
Falcon-512 & 1,622~B & 108.6~M gas & $\sim$2$\times$ current Ethereum config. \\
ML-DSA-44 & 3,790~B & 46.5~M gas & Below 60~M target \\
\bottomrule
\end{tabular}
\end{table}
We observe that ECDSA transactions remain far below the data-limited regime under any realistic gas configuration.
Falcon-512 approaches data saturation only at approximately twice the current Ethereum gas limits (45~M~\cite{eth_gas_limit_45M_2025}).
In contrast, ML-DSA-44 becomes data-limited within the range of existing Ethereum target parameters (60~M~\cite{eth_gas_limit_60M_2025}), indicating that block size, not gas or CPU, becomes the constraint.
\footnotetext{The results for CATX ECDSA and legacy transactions are reported to be equivalent.}

Next, we present the resulting throughput measurements in Table~\ref{tab:throughput}, allowing us to identify the dominant bottleneck for each transaction type.
% \vspace{-10pt}
\begin{table}[htbp]
\centering
\caption{Transactions per second (TPS), gas and block size consumption, and the primary performance bottleneck for legacy \texttt{0x02} and CATX transaction types.}
\label{tab:throughput}
\setlength{\tabcolsep}{5pt}
\begin{tabular}{@{}lrrrrrr@{}}
\toprule
\textbf{Type} & \textbf{Peak TPS} & \textbf{Sustained TPS} & \textbf{Gas\%} & \textbf{Size\%} & \textbf{Bottleneck} \\
\midrule
Legacy & 4,021 & 3,963 & 19\% & 1\% & CPU \\
ECDSA & 4,032 & 3,964 & 22\% & 2\% & CPU \\
Falcon-512 & 3,730 & 1,625~B & 78\% & 72\% & Mixed \\
ML-DSA-44 & 2,210 & 3,794~B & 46\% & \textbf{100\%} & \textbf{Block size} \\
\bottomrule
\end{tabular}%
\end{table}
% \vspace{-10pt}
% Figure~\ref{fig:tps-comparison} visualizes these results.
The CATX transaction format introduces no measurable overhead: ECDSA achieves the same throughput as legacy EIP-1559 transactions, with both ultimately limited by the slot time. 
Falcon-512 maintains high throughput despite significantly larger signatures, operating near the transition between gas- and size-limited regimes. 
In contrast, ML-DSA-44 exhibits lower throughput because blocks reach the 8~MB size limit before exhausting the available gas.

This behavior arises from the fact that overall transactions per second (TPS) is directly determined by transactions per block (TPB) and blocks per second (BPS): $\text{TPS}=\text{TPB}\times\text{BPS}$. 
At the extreme, although ML-DSA-44 accommodates more transactions per block due to larger block sizes, the slower block production rate dominates, yielding lower net TPS. 
Figure~\ref{fig:block-size-impact} highlights the negative correlation between block size and block production rate across all evaluated signature schemes. 
Hence, faster cryptographic verification does not translate into higher system throughput once block propagation becomes the limiting factor.

For completeness, exploratory end-to-end runs with SLH-DSA showed that SHAKE-128f reached only around 200~TPS, whereas SHAKE-128s was not practical to benchmark because signing is 8{,}576$\times$ slower than ECDSA.
% Requires: \usepackage{pgfplots}
%          \usepgfplotslibrary{fillbetween}
%          \pgfplotsset{compat=1.18}
%

\begin{figure}[!t]
    \centering
    \begin{tikzpicture}
    \begin{axis}[   
        width=\textwidth,
        height=4.5cm,
        xmode=log,
        xlabel={Block size (MB)},
        ylabel={Blocks per second (BPS)},
        title={Block size impact on BPS},
        grid=both,
        grid style={dashed, gray!30},
        % legend style={at={(0.3,0.05)}, anchor=south east, font=\small}, % below
        legend style={at={(0.81,0.63)}, anchor=south east, font=\tiny}, % above
        % log ticks with fixed point,
        xmin=0.005,
        xmax=37,
        ymin=0.5,
        ymax=6.5,
        ytick={1,2,3,4,5,6},
        % minor y tick num=1,
        % major grid style={dashed, gray!30},
        % minor grid style={dotted, gray!20},
        % yminorgrids=true,
    ]
    
    % Plot order: scatter (background) -> bands -> lines (foreground)
    
    % ===== ECDSA =====
    % Scatter points (raw data) - plotted first so they're in background
    \addplot[
        only marks,
        mark=*,
        mark size=0.35pt,
        color=ecdsa_line,
        opacity=0.05,
        forget plot
    ] table[x=blockSizeMb, y=bps, col sep=comma] {figures/data/block_size_bps_ecdsa_raw.csv};
    
    % Error band - upper bound
    \addplot[
        draw=none,
        name path=ecdsa_upper,
        forget plot
    ] table[x=blockSizeMb, y=stdUpper, col sep=comma] {figures/data/block_size_bps_ecdsa.csv};
    
    % Error band - lower bound
    \addplot[
        draw=none,
        name path=ecdsa_lower,
        forget plot
    ] table[x=blockSizeMb, y=stdLower, col sep=comma] {figures/data/block_size_bps_ecdsa.csv};
    
    % Error band - fill between
    \addplot[
        fill=ecdsa_band,
        fill opacity=0.4,
        forget plot
    ] fill between[of=ecdsa_upper and ecdsa_lower];
    
    % Mean line - plotted last so it's on top
    \addplot[
        color=ecdsa_line,
        line width=0.5pt,
        mark=none
    ] table[x=blockSizeMb, y=meanBps, col sep=comma] {figures/data/block_size_bps_ecdsa.csv};
    \addlegendentry{ECDSA}
    
    % ===== Falcon-512 =====
    % Scatter points (raw data)
    \addplot[
        only marks,
        mark=*,
        mark size=0.35pt,
        color=falcon_line,
        opacity=0.05,
        forget plot
    ] table[x=blockSizeMb, y=bps, col sep=comma] {figures/data/block_size_bps_falcon_512_raw.csv};
    
    % Error band - upper bound
    \addplot[
        draw=none,
        name path=falcon_upper,
        forget plot
    ] table[x=blockSizeMb, y=stdUpper, col sep=comma] {figures/data/block_size_bps_falcon_512.csv};
    
    % Error band - lower bound
    \addplot[
        draw=none,
        name path=falcon_lower,
        forget plot
    ] table[x=blockSizeMb, y=stdLower, col sep=comma] {figures/data/block_size_bps_falcon_512.csv};
    
    % Error band - fill between
    \addplot[
        fill=falcon_band,
        fill opacity=0.4,
        forget plot
    ] fill between[of=falcon_upper and falcon_lower];
    
    % Mean line
    \addplot[
        color=falcon_line,
        line width=0.5pt,
        mark=none
    ] table[x=blockSizeMb, y=meanBps, col sep=comma] {figures/data/block_size_bps_falcon_512.csv};
    \addlegendentry{Falcon-512}
    
    % ===== ML-DSA-44 =====
    % Scatter points (raw data)
    \addplot[
        only marks,
        mark=*,
        mark size=0.35pt,
        color=mldsa_line,
        opacity=0.05,
        forget plot
    ] table[x=blockSizeMb, y=bps, col sep=comma] {figures/data/block_size_bps_ml_dsa_44_raw.csv};
    
    % Error band - upper bound
    \addplot[
        draw=none,
        name path=mldsa_upper,
        forget plot
    ] table[x=blockSizeMb, y=stdUpper, col sep=comma] {figures/data/block_size_bps_ml_dsa_44.csv};
    
    % Error band - lower bound
    \addplot[
        draw=none,
        name path=mldsa_lower,
        forget plot
    ] table[x=blockSizeMb, y=stdLower, col sep=comma] {figures/data/block_size_bps_ml_dsa_44.csv};
    
    % Error band - fill between
    \addplot[
        fill=mldsa_band,
        fill opacity=0.4,
        forget plot
    ] fill between[of=mldsa_upper and mldsa_lower];
    
    % Mean line
    \addplot[
        color=mldsa_line,
        line width=0.5pt,
        mark=none
    ] table[x=blockSizeMb, y=meanBps, col sep=comma] {figures/data/block_size_bps_ml_dsa_44.csv};
    \addlegendentry{ML-DSA-44}

    % 8 MB block limit vertical line
    \addplot[dashed, thick, gray] coordinates {(8, 0) (8, 7)};
    \node[font=\scriptsize, gray, anchor=west] at (axis cs:8.2, 3.5) {8 MB limit};
        
    \end{axis}
    \end{tikzpicture}
    \caption{Block size impact on BPS. Error bands show ±1 standard deviation. Based on observation from a live blockchain node (aggregated into 15-second windows), showing BPS behavior when each scheme dominates block production (ECDSA: 2,294 windows; Falcon-512: 1,104 windows; ML-DSA-44: 493 windows)}
    \label{fig:block-size-impact}
\end{figure}
    